\section{Analysis of Model-\texorpdfstring{$X$}{X} (\texorpdfstring{$\gamma_1,\gamma_2>0,\ \theta_1,\theta_2>0$}{gamma1,gamma2>0, theta1,theta2>0})}\label{sec:model_x_anal}

Model-$X$ is the full model: it retains \emph{both} bidirectional jockeying
($\gamma_1,\gamma_2>0$) \emph{and} two-class abandonment ($\theta_1,\theta_2>0$). It is
the parent of the hierarchy, and every other model analysed in this work is recovered
from it by zeroing parameters. Jockeying conserves the total number of customers while
abandonment removes them; the second mechanism makes the chain positive recurrent for
all loads, so the stationary distribution always exists. As in Model-$B$, jockeying
breaks the per-class Poisson/renewal structure~---~a Class-2 customer can no longer
identify its effective service time with a Class-1 busy period~---~and abandonment,
although it leaves each arrival stream Poisson, couples the two queues through the
state-dependent rates $\theta_i n_i$; in either case the Pollaczek--Khinchine
probabilistic route is unavailable. Applying the method of characteristics to the
fundamental equation still yields a general solution, but closing it~---~i.e.,
determining the unknown boundary function $P_y(y)$~---~requires solving a Volterra-type
integral equation that does not admit a closed form. In fact the obstruction is
strictly more severe than in Model-$B$, on two counts developed below: the
characteristics are no longer straight lines but integral curves of a nonlinear ODE,
and the diagonal value $P(z,z)$ that anchored the reduction in Model-$B$ is itself no
longer known. Model-$X$ is therefore intractable and is left open; its role is to
expose the precise obstructions that its tractable specialisations
$A$, $B$, $B_2$ and $C_2$ each circumvent. The queue layout is the left panel of
Figure~\ref{fig:models-queue-overview}.

\textbf{Stability.} Model-$X$ has abandonment from both classes, so the total departure
rate in state $(n_1,n_2)$ is $\mu+\theta_1 n_1+\theta_2 n_2$, which grows without bound
in every direction of the $(n_1,n_2)$-plane. The chain is therefore positive recurrent
for \emph{all} $\lambda_1,\lambda_2,\mu>0$ and all $\theta_1,\theta_2>0$; no
traffic-intensity restriction is required (cf.\ the Erlang-A queue,
\S\ref{ssec:erlangA}). Setting $\theta_1=\theta_2=0$ removes this and returns the
Model-$B$/Model-$A$ condition $\rho=(\lambda_1+\lambda_2)/\mu<1$.

\begin{corollary}\label{cor:X:pi00}
    The empty-system probability $\pi_0$ (server idle, no one waiting) and the
    both-queues-empty busy probability $\pi(0,0)$ (server busy, $N_1=N_2=0$) of
    Model-$X$ satisfy the universal idle balance
    \begin{equation}
        \pi(0,0) = \rho\,\pi_0,
        \qquad \rho=\frac{\lambda_1+\lambda_2}{\mu},
        \label{eq:X:pi00}
    \end{equation}
    as in every model of this work~\eqref{eq:gen:pi00_general}. Unlike Models~$A$
    and~$B$, however, abandonment makes the total-customer process a state-dependent
    birth--death chain rather than an $M/M/1$, so $\pi_0\neq 1-\rho$: the idle
    probability depends on $\theta_1,\theta_2$ and is fixed only by the global
    normalisation $P(1,1)=1-\pi_0$ once $P$ is known. Correspondingly, the diagonal of
    the PGF no longer collapses to $\pi(0,0)/(1-\rho z)$; setting $x=y=z$ in the
    fundamental equation leaves the abandonment-corrected relation
    \begin{equation}
        (\mu-\lambda z)\,P(z,z)
        + z\bigl[\theta_1\,\partial_x P + \theta_2\,\partial_y P\bigr]_{(z,z)}
        = \mu\,\pi(0,0),
        \qquad \lambda=\lambda_1+\lambda_2.
        \label{eq:X:diagonal}
    \end{equation}
    As $\theta_1,\theta_2\to 0^+$ the gradient terms drop out and~\eqref{eq:X:diagonal}
    returns the Model-$B$ identities $P(z,z)=\pi(0,0)/(1-\rho z)$, $\pi_0=1-\rho$ and
    $\pi(0,0)=\rho(1-\rho)$~\eqref{eq:B:pi00},~\eqref{eq:A:pi0_pi00}.
\end{corollary}

\begin{proof}
    Setting $x=y$ in the fundamental
    equation~\eqref{eq:gen:fundamental_equation}, both jockeying terms carry the factor
    $(x-y)$ and vanish, as does the boundary term $\mu(x-y)P_y(y)$; the abandonment
    terms, carrying instead the factor $(x-1)$, survive. The coefficient of $P$ becomes
    $(\lambda_1+\lambda_2+\mu)z^2-\mu z-(\lambda_1+\lambda_2)z^3 = z(z-1)(\mu-\lambda z)$,
    and the right-hand side becomes $-\mu z(1-z)\pi(0,0)=\mu z(z-1)\pi(0,0)$; dividing
    the surviving terms by the common factor $z(z-1)$ (for $z\in(0,1)$) leaves
    \begin{align}
        z(z-1)(\mu-\lambda z)P(z,z)
        + z^2(z-1)\bigl[\theta_1\,\partial_x P + \theta_2\,\partial_y P\bigr]_{(z,z)}
        = \mu z(z-1)\,\pi(0,0),
        \label{eq:X:diagonal_full}
    \end{align}
    which is~\eqref{eq:X:diagonal}. The idle state $(0)$ exchanges probability flux only
    with the both-queues-empty busy state, so global balance across the idle--busy cut
    reads $\lambda\pi_0=\mu\pi(0,0)$~\eqref{eq:gen:balance_idle} (an arrival at total rate
    $\lambda$ leaves $(0)$; a service completion at rate $\mu$ returns to it), whence
    $\pi(0,0)=\rho\pi_0$, which is~\eqref{eq:X:pi00}. In Models~$A$ and~$B$ the gradient
    terms in~\eqref{eq:X:diagonal} are absent ($\theta_1=\theta_2=0$), so
    $P(z,z)=\pi(0,0)/(1-\rho z)$; evaluating at $z=1$ with $P(1,1)=1-\pi_0$ then forces
    $(1-\pi_0)(1-\rho)=\rho\pi_0$, i.e.\ $\pi_0=1-\rho$. With abandonment those gradient
    terms are present and \emph{not} known~---~exactly the phenomenon isolated for
    Model-$C_2$~\eqref{eq:C2:diag}~---~so the diagonal no longer determines $P(z,z)$, and
    $\pi_0$ cannot be read off in closed form: it follows only from the full solution
    via normalisation.
\end{proof}

\begin{claim}\label{clm:X:volterra}
    Under positive recurrence, the bivariate PGF of Model-$X$ satisfies, on each
    characteristic curve $\Phi(x,y)=C_1$ of the fundamental equation, the integral
    representation
    \begin{equation}
        P(x,y(x)) \;=\; P_h(x,y(x))\int_{0}^{x}
        \frac{h\bigl(t;\,C_1\bigr)}{P_h\bigl(t,\,y(t)\bigr)}\,dt,
        \label{eq:X:integral_rep}
    \end{equation}
    where $P_h$ is the homogeneous solution~\eqref{eq:X:Ph}, $y(t)=y(t;C_1)$ is the
    characteristic parametrisation, $\Phi$ is the first integral~\eqref{eq:X:first_integral}
    (a transcendental function of $(x,y)$, not the linear form of Model-$B$), and
    $h(t;C_1)$ is a forcing function depending on the unknown boundary trace
    $P_y(y)=P(0,y)$. Evaluating~\eqref{eq:X:integral_rep} along the characteristic
    through a diagonal point $(z,z)$ reduces the problem to a \emph{Volterra integral
    equation of the first kind} for $P_y$:
    \begin{equation}
        \mathcal{F}(z) \;=\; \int_0^z \mathcal{K}(z,t)\,P_y\!\bigl(y(t;z)\bigr)\,dt,
        \label{eq:X:volterra}
    \end{equation}
    whose kernel $\mathcal{K}$ and right-hand side $\mathcal{F}$ are defined
    in~\eqref{eq:X:kernel}--\eqref{eq:X:rhs} below. The kernel $\mathcal{K}$ is built
    from the non-elementary homogeneous solution $P_h$ and does not simplify for general
    $(\gamma_1,\gamma_2,\theta_1,\theta_2)$; moreover, in contrast to Model-$B$, the
    right-hand side $\mathcal{F}$ is \emph{not} fully known~---~it is coupled
    through~\eqref{eq:X:diagonal} to the diagonal gradient
    $[\theta_1\partial_x P+\theta_2\partial_y P]_{(z,z)}$. Equation~\eqref{eq:X:volterra}
    therefore does not reduce to a closed-form solution for $P_y$, and Model-$X$ is left
    open.
\end{claim}

\begin{proof}[Proof of Claim~\ref{clm:X:volterra}]
The proof follows the five steps of Model-$B$: (i)~homogeneous solution via
characteristics, (ii)~behaviour of $P_h$ at the boundaries, (iii)~particular solution
via variation of parameters, (iv)~pinning the free function from the boundary condition,
and (v)~diagonal evaluation yielding the Volterra equation. At each step we point out
where the abandonment terms $\theta_i$ obstruct the closed-form algebra available in
Model-$B$.

% ─────────────────────────────────────────────────────────────────────────────
\medskip\noindent\textbf{Step~1: Homogeneous solution via the method of characteristics.}

The fundamental equation~\eqref{eq:gen:fundamental_equation} is a first-order linear PDE
in $P(x,y)$. Setting the right-hand side to zero and dividing through by $y$ yields the
homogeneous PDE
\begin{equation}
    x\bigl[\gamma_1(x-y)+\theta_1(x-1)\bigr]\,\frac{\partial P}{\partial x}
    + x\bigl[\gamma_2(y-x)+\theta_2(y-1)\bigr]\,\frac{\partial P}{\partial y}
    = -A(x,y)\,P,
    \label{eq:X:homogeneous_pde}
\end{equation}
with the same coefficient polynomial as Model-$B$,
\begin{equation}
    A(x,y) := (\lambda_1+\lambda_2+\mu)x - \mu - \lambda_1 x^2 - \lambda_2 xy ;
    \label{eq:X:Adef}
\end{equation}
abandonment changes only the two convection (drift) coefficients, adding the terms
$\theta_1(x-1)$ and $\theta_2(y-1)$. Following the method of characteristics
(\S\ref{sssec:moc}), we associate to~\eqref{eq:X:homogeneous_pde} the characteristic
system
\begin{equation}
    \frac{dx}{x\bigl[\gamma_1(x-y)+\theta_1(x-1)\bigr]}
    = \frac{dy}{x\bigl[\gamma_2(y-x)+\theta_2(y-1)\bigr]}
    = \frac{dP}{-A(x,y)\,P}.
    \label{eq:X:char_system}
\end{equation}
Taking the ratio of the first two differentials, the common factor $x$ cancels:
\begin{equation}
    \frac{dy}{dx}
    = \frac{\gamma_2(y-x)+\theta_2(y-1)}{\gamma_1(x-y)+\theta_1(x-1)}
    = \frac{-\gamma_2 x + (\gamma_2+\theta_2)y - \theta_2}
           {(\gamma_1+\theta_1)x - \gamma_1 y - \theta_1}.
    \label{eq:X:char_ode}
\end{equation}
In Model-$B$ ($\theta_1=\theta_2=0$) both affine forms shared the factor $(x-y)$ and the
ratio collapsed to the constant $-\gamma_2/\gamma_1$, integrating to the straight lines
$\gamma_2 x+\gamma_1 y=\mathrm{const}$. With abandonment this collapse no longer occurs:
\eqref{eq:X:char_ode} is a genuine first-order ODE whose right-hand side is a ratio of
two affine forms, and its integral curves are not straight lines.

Both affine forms in~\eqref{eq:X:char_ode} vanish simultaneously only at the single
point $(x,y)=(1,1)$ (the unique intersection of the lines
$(\gamma_1+\theta_1)x-\gamma_1 y-\theta_1=0$ and
$-\gamma_2 x+(\gamma_2+\theta_2)y-\theta_2=0$), which is therefore the lone singular
point of the characteristic field. Shifting to it via $X=x-1,\ Y=y-1$ removes the
constant terms and renders~\eqref{eq:X:char_ode} homogeneous:
\begin{equation}
    \frac{dY}{dX}
    = \frac{-\gamma_2 X + (\gamma_2+\theta_2)Y}{(\gamma_1+\theta_1)X - \gamma_1 Y}.
    \label{eq:X:char_shifted}
\end{equation}
The substitution $v=Y/X$ then separates the variables,
\begin{equation}
    \frac{(\gamma_1+\theta_1)-\gamma_1 v}{Q(v)}\,dv = \frac{dX}{X},
    \qquad
    Q(v) := \gamma_1 v^2 + (\gamma_2+\theta_2-\gamma_1-\theta_1)\,v - \gamma_2 .
    \label{eq:X:char_separated}
\end{equation}
The quadratic $Q$ has discriminant
$D=(\gamma_2+\theta_2-\gamma_1-\theta_1)^2+4\gamma_1\gamma_2>0$ and real roots of opposite
sign (their product is $-\gamma_2/\gamma_1<0$),
\begin{equation}
    v_{\pm}=\frac{-(\gamma_2+\theta_2-\gamma_1-\theta_1)\pm\sqrt{D}}{2\gamma_1},
    \qquad v_-<0<v_+ .
    \label{eq:X:roots}
\end{equation}
Decomposing the left-hand side of~\eqref{eq:X:char_separated} by partial fractions over
$Q(v)=\gamma_1(v-v_+)(v-v_-)$,
\begin{equation}
    \frac{(\gamma_1+\theta_1)-\gamma_1 v}{\gamma_1(v-v_+)(v-v_-)}
    = \frac{a_+}{v-v_+} + \frac{a_-}{v-v_-},
    \qquad
    a_{\pm}=\frac{(\gamma_1+\theta_1)-\gamma_1 v_{\pm}}{\gamma_1\,(v_{\pm}-v_{\mp})},
    \qquad a_++a_-=-1,
    \label{eq:X:apm}
\end{equation}
(the relation $a_++a_-=-1$ is the residue of the integrand at $v=\infty$), and
integrating gives $a_+\ln(v-v_+)+a_-\ln(v-v_-)=\ln X+\mathrm{const}$. Exponentiating and
returning to $(x,y)$ via $v-v_\pm=(Y-v_\pm X)/X$ and $a_++a_-=-1$, the powers of $X$
cancel and we obtain the first integral
\begin{equation}
    \Phi(x,y)
    := \bigl[(y-1)-v_+(x-1)\bigr]^{a_+}\,\bigl[(y-1)-v_-(x-1)\bigr]^{a_-}
    = C_1,
    \label{eq:X:first_integral}
\end{equation}
constant along each characteristic. Unlike Model-$B$'s linear constant
$\gamma_2 x+\gamma_1 y$, the first integral~\eqref{eq:X:first_integral} is a product of
non-integer powers of two linear forms~---~a transcendental level curve. As a
consistency check, letting $\theta_1,\theta_2\to0^+$ gives
$Q(v)\to(\gamma_1 v+\gamma_2)(v-1)$, so $v_+\to1$, $v_-\to-\gamma_2/\gamma_1$, and
$(a_+,a_-)\to(0,-1)$; then
$\Phi(x,y)\to\bigl[(y-1)+(\gamma_2/\gamma_1)(x-1)\bigr]^{-1}$, whose level curves are
exactly the straight lines $\gamma_2 x+\gamma_1 y=\mathrm{const}$ of Model-$B$.

To find $P$ along a characteristic we take the ratio of the first and third differentials
in~\eqref{eq:X:char_system},
\begin{equation}
    \frac{dP}{P}
    = \frac{-A\bigl(x,y(x)\bigr)}
           {x\bigl[\gamma_1(x-y(x))+\theta_1(x-1)\bigr]}\,dx ,
    \label{eq:X:dPP}
\end{equation}
where $y(x)=y(x;C_1)$ is the characteristic~\eqref{eq:X:first_integral}. Integrating
yields the general homogeneous solution
\begin{equation}
    P_h(x,y)
    = F(C_1)\,\exp\!\left(
        -\int_{x_\star}^{x}
        \frac{A\bigl(s,y(s)\bigr)}
             {s\bigl[\gamma_1(s-y(s))+\theta_1(s-1)\bigr]}\,ds
    \right),
    \qquad C_1=\Phi(x,y),
    \label{eq:X:Ph}
\end{equation}
where $F$ is an arbitrary differentiable function of the characteristic constant and
$x_\star$ is an arbitrary base point (its choice is absorbed into $F$). Here the
parallel with Model-$B$ ends in closed form: because $y(s)$ is the transcendental
characteristic~\eqref{eq:X:first_integral}, the integrand in~\eqref{eq:X:Ph} is not a
rational function of $s$, and the quadrature is not elementary. Model-$B$, by contrast,
could substitute the linear $y(s)$, partial-fraction the rational integrand, and read
off the explicit factors $x^{-\mu/C_1}(x-y)^{E}$ with a closed-form exponent $E$. No such
explicit $P_h$ exists for Model-$X$.

% ─────────────────────────────────────────────────────────────────────────────
\medskip\noindent\textbf{Step~2: Boundary behaviour of $P_h$.}

Although the quadrature~\eqref{eq:X:Ph} is not elementary, the two qualitative features
needed downstream survive, and a third feature of Model-$B$ is lost.

\smallskip\noindent\textit{(a) Blow-up at $x=0$.}
A characteristic meets the axis $x=0$ at a finite height $y_0:=y(0;C_1)\in(0,1]$, where
$A(0,y_0)=-\mu$ and the convection coefficient is
$\gamma_1(0-y_0)+\theta_1(0-1)=-(\gamma_1 y_0+\theta_1)<0$. Hence the integrand
in~\eqref{eq:X:Ph} behaves like $-\mu/[(\gamma_1 y_0+\theta_1)\,s]$ as $s\to0^+$, so
\begin{equation}
    P_h\bigl(x,y(x)\bigr)\ \sim\ x^{-\mu/(\gamma_1 y_0+\theta_1)}\ \longrightarrow\ +\infty
    \qquad (x\to0^+),
    \label{eq:X:Ph_at_zero}
\end{equation}
with a strictly negative exponent because $\gamma_1 y_0+\theta_1>0$. This is exactly the
factor $x^{-\mu/C_1}$ of Model-$B$, now with the denominator shifted by $\theta_1$; it
drives the pinning of $F$ in Step~4.

\smallskip\noindent\textit{(b) Regularity on the diagonal.}
In Model-$B$ the convection coefficient was $\gamma_1 x(x-y)$, which \emph{vanishes} on
the diagonal $x=y$; the resulting factor $(x-y)^{E}$ with $E>0$ made $P_h\to0$ there and
forced the delicate Hadamard finite-part regularisation of Model-$B$'s Step~5. For
Model-$X$ the coefficient is $x[\gamma_1(x-y)+\theta_1(x-1)]$, which at $x=y=z$ equals
$z\,\theta_1(z-1)\neq0$ for $z\in(0,1)$: the abandonment term $\theta_1(x-1)$ keeps the
characteristic field away from zero on the diagonal. Consequently the integrand
in~\eqref{eq:X:Ph} is regular as $x\to z$ along the diagonal-crossing characteristic, and
\begin{equation}
    P_h^{\,\mathrm{diag}}(z):=\lim_{x\to z}P_h\bigl(x,y(x;z)\bigr)
    \quad\text{is finite and non-zero.}
    \label{eq:X:Ph_diag}
\end{equation}
The diagonal singularity of Model-$B$ thus disappears; as we will see in Step~5, the
obstruction migrates instead to the right-hand side of the Volterra equation.

\smallskip\noindent\textit{Branch convention.}
As in Model-$B$, the non-integer powers in~\eqref{eq:X:first_integral} and the quadrature
in~\eqref{eq:X:Ph} may carry a complex phase off the real region; this phase is absorbed
into the arbitrary function $F(C_1)$ and cancels exactly in the ratio
$P_h(x,y(x))/P_h(t,y(t))$ appearing in~\eqref{eq:X:integral_rep}, since numerator and
denominator lie on the same characteristic (same $C_1$).

% ─────────────────────────────────────────────────────────────────────────────
\medskip\noindent\textbf{Step~3: Particular solution via variation of parameters.}

We now restrict the full (inhomogeneous) fundamental
equation~\eqref{eq:gen:fundamental_equation} to a fixed characteristic. Along
$y=y(x;C_1)$ the total derivative is
$\frac{d}{dx}P(x,y(x))=\partial_x P + y'(x)\,\partial_y P$, and by the characteristic
ODE~\eqref{eq:X:char_ode} the bracket multiplying $\partial_y P$ equals
$[\gamma_1(x-y)+\theta_1(x-1)]\,y'(x)$, so the two convection terms combine:
\[
    xy\bigl[\gamma_1(x-y)+\theta_1(x-1)\bigr]\partial_x P
    + xy\bigl[\gamma_2(y-x)+\theta_2(y-1)\bigr]\partial_y P
    = xy\bigl[\gamma_1(x-y)+\theta_1(x-1)\bigr]\frac{d}{dx}P(x,y(x)).
\]
Substituting into~\eqref{eq:gen:fundamental_equation} and dividing by
$xy\bigl[\gamma_1(x-y)+\theta_1(x-1)\bigr]$ reduces the PDE to the first-order linear ODE
(for fixed $C_1$)
\begin{equation}
    \frac{dP}{dx} = f(x;C_1)\,P + h(x;C_1),
    \label{eq:X:ode_char}
\end{equation}
with coefficient and forcing
\begin{align}
    f(x;C_1) &= -\frac{A\bigl(x,y(x)\bigr)}
                       {x\bigl[\gamma_1(x-y(x))+\theta_1(x-1)\bigr]},
    \label{eq:X:fdef}\\[4pt]
    h(x;C_1) &= \frac{G\bigl(x,y(x)\bigr)}
                     {x\,y(x)\bigl[\gamma_1(x-y(x))+\theta_1(x-1)\bigr]},
    \label{eq:X:hdef}
\end{align}
where $A$ is as in~\eqref{eq:X:Adef} and the right-hand side of the fundamental equation
is, as in Model-$B$,
\begin{equation}
    G(x,y) := \mu(x-y)\,P_y(y) - \mu x(1-y)\,\pi(0,0).
    \label{eq:X:Gdef}
\end{equation}
Comparing~\eqref{eq:X:fdef}--\eqref{eq:X:hdef} with their Model-$B$ counterparts, the
denominator $\gamma_1 x(x-y)$ is replaced throughout by
$x[\gamma_1(x-y)+\theta_1(x-1)]$~---~the only footprint of abandonment, and the source of
all the difficulty. Since $P_h$ from~\eqref{eq:X:Ph} satisfies the homogeneous equation
$\frac{d}{dx}P_h=f(x;C_1)P_h$, it is the fundamental solution
of~\eqref{eq:X:ode_char}; applying variation of parameters (\S\ref{sssec:vop}), the
general solution of the inhomogeneous ODE is
\begin{equation}
    P(x,y(x)) = P_h(x,y(x))\!\left[F(C_1)
        + \int_{x_0}^{x} \frac{h(t;C_1)}{P_h(t,y(t))}\,dt\right].
    \label{eq:X:variation}
\end{equation}

% ─────────────────────────────────────────────────────────────────────────────
\medskip\noindent\textbf{Step~4: Pinning the free function from the boundary condition at $x=0$.}

Along any characteristic the endpoint at $x=0$ is $y(0)=y_0$, giving
$P(0,y_0)=P_y(y_0)$, which is finite. However, by~\eqref{eq:X:Ph_at_zero} the homogeneous
factor blows up, $P_h(x,y(x))\to+\infty$ as $x\to0^+$, while the integral
in~\eqref{eq:X:variation} vanishes as $x_0\to0$ (the integration interval collapses).
Therefore
\[
    \lim_{x\to0^+}P(x,y(x))
    = \lim_{x\to0^+}P_h(x,y(x))\cdot F(C_1)
    = +\infty\cdot F(C_1),
\]
which can be finite only if $F(C_1)=0$. Setting $x_0=0$ and $F(C_1)=0$
in~\eqref{eq:X:variation} gives the integral representation of
Claim~\ref{clm:X:volterra},
\begin{equation}
    P(x,y(x)) = P_h(x,y(x))\int_{0}^{x} \frac{h(t;C_1)}{P_h(t,y(t))}\,dt.
    \label{eq:X:integral_rep_body}
\end{equation}
This argument is identical to Model-$B$'s; abandonment shifts the blow-up exponent
in~\eqref{eq:X:Ph_at_zero} from $\mu/C_1$ to $\mu/(\gamma_1 y_0+\theta_1)$ but leaves it
strictly positive, so the conclusion $F\equiv0$ is unchanged.

% ─────────────────────────────────────────────────────────────────────────────
\medskip\noindent\textbf{Step~5: Reduction to the Volterra equation.}

We simplify the forcing $h$ on the characteristic. Substituting $G$
from~\eqref{eq:X:Gdef} into~\eqref{eq:X:hdef} and cancelling one factor of $t$ in the
second term,
\begin{equation}
    h(t;C_1)
    = \frac{\mu\,(t-y(t))}{t\,y(t)\bigl[\gamma_1(t-y(t))+\theta_1(t-1)\bigr]}\,P_y\!\bigl(y(t)\bigr)
    - \frac{\mu\,(1-y(t))}{y(t)\bigl[\gamma_1(t-y(t))+\theta_1(t-1)\bigr]}\,\pi(0,0).
    \label{eq:X:h_simplified}
\end{equation}
We evaluate~\eqref{eq:X:integral_rep_body} along the characteristic through a diagonal
point $(z,z)$, that is, the curve with $C_1=\Phi(z,z)$; its parametrisation $y(t;z)$
satisfies $y(z;z)=z$. Setting $x=z$,
\begin{equation}
    P(z,z) = P_h(z,z)\int_0^z \frac{h\bigl(t;\Phi(z,z)\bigr)}{P_h\bigl(t,y(t;z)\bigr)}\,dt.
    \label{eq:X:diag_eval}
\end{equation}
By the regularity~\eqref{eq:X:Ph_diag}, $P_h(z,z)=P_h^{\,\mathrm{diag}}(z)$ is finite and
non-zero and the integrand is regular as $t\to z$ (its numerator
$t-y(t;z)\to0$), so~\eqref{eq:X:diag_eval} is a \emph{proper} integral~---~no
finite-part regularisation is needed, in contrast to Model-$B$. Substituting the two
pieces of~\eqref{eq:X:h_simplified} and separating the term containing the unknown $P_y$
from the known remainder, define the kernel and remainder
\begin{align}
    \mathcal{K}(z,t) &:= \frac{\mu\,\bigl(t-y(t;z)\bigr)\,P_h^{\,\mathrm{diag}}(z)}
        {t\,y(t;z)\bigl[\gamma_1(t-y(t;z))+\theta_1(t-1)\bigr]\,P_h\bigl(t,y(t;z)\bigr)},
    \label{eq:X:kernel}\\[6pt]
    \mathcal{R}(z) &:= -\,\mu\,\pi(0,0)\,P_h^{\,\mathrm{diag}}(z)
        \int_0^z \frac{1-y(t;z)}
        {y(t;z)\bigl[\gamma_1(t-y(t;z))+\theta_1(t-1)\bigr]\,P_h\bigl(t,y(t;z)\bigr)}\,dt,
    \label{eq:X:Rdef}
\end{align}
so that~\eqref{eq:X:diag_eval} reads
$P(z,z)=\int_0^z\mathcal{K}(z,t)\,P_y(y(t;z))\,dt+\mathcal{R}(z)$. Setting
\begin{equation}
    \mathcal{F}(z) := P(z,z) - \mathcal{R}(z)
    = \underbrace{\frac{\pi(0,0)}{1-\rho z}}_{\text{Model-}B\text{ value}}
      \;-\;\underbrace{\frac{z}{\mu-\lambda z}
        \bigl[\theta_1\partial_x P + \theta_2\partial_y P\bigr]_{(z,z)}}_{\text{abandonment coupling}}
      \;-\;\mathcal{R}(z),
    \label{eq:X:rhs}
\end{equation}
where the expression for $P(z,z)$ is the abandonment-corrected
diagonal~\eqref{eq:X:diagonal}, equation~\eqref{eq:X:diag_eval} becomes the Volterra
integral equation of the first kind~\eqref{eq:X:volterra}, confirming
Claim~\ref{clm:X:volterra}.
% [AUTHOR: Model-$B$ closes because, with theta_i=0, (i) the characteristics are straight
% lines so K is explicit, and (ii) the diagonal gradient term in F vanishes, leaving
% F(z)=pi(0,0)/(1-rho z) - R(z) fully known -- a clean (if unsolved) first-kind Volterra
% equation. Model-$X$ obstructs BOTH: K rests on the non-elementary P_h, AND F is coupled
% to the unknown diagonal gradient [theta_1 P_x + theta_2 P_y]_(z,z). The second coupling
% is the genuinely new obstruction; it is the same phenomenon that prevents the diagonal
% from closing in Model C_2 (eq:C2:diag). Flagged so this reads as a deliberate open end.]
Two features distinguish~\eqref{eq:X:volterra} from its Model-$B$ analogue, and both make
Model-$X$ strictly less tractable. First, the kernel~\eqref{eq:X:kernel} is built from the
non-elementary homogeneous solution $P_h$ along transcendental characteristics, so it has
no closed form for general parameters (in Model-$B$, $P_h$ and the straight-line
characteristics were explicit). Second, and decisively, the right-hand
side~\eqref{eq:X:rhs} is \emph{not} known data: through~\eqref{eq:X:diagonal} it is
coupled to the diagonal gradient $[\theta_1\partial_x P+\theta_2\partial_y P]_{(z,z)}$,
which is precisely the quantity that abandonment prevents the diagonal from
determining~\eqref{eq:C2:diag}. Where Model-$B$ left a clean first-kind Volterra equation
with explicit kernel and known forcing (itself already beyond closed-form inversion),
Model-$X$ leaves one whose kernel is implicit and whose forcing is entangled with the
unknown solution. Resolving it lies well beyond the scope of this thesis.
\end{proof}

Model-$X$ is therefore left open at this point. Its role in the hierarchy is to make the
two obstructions explicit and to show how each tractable specialisation removes them. The
limit $\theta_1=\theta_2\to0$ (Model-$B$) restores both at once: the
characteristics~\eqref{eq:X:first_integral} straighten to the lines
$\gamma_2 x+\gamma_1 y=\mathrm{const}$, so the kernel~\eqref{eq:X:kernel} becomes
explicit, and the gradient coupling in~\eqref{eq:X:rhs} vanishes, so
$\mathcal{F}(z)=\pi(0,0)/(1-\rho z)-\mathcal{R}(z)$ becomes known data~---~leaving the
clean (still unsolved) Volterra equation of Section~\ref{sec:model_b}. Removing one
mechanism at a time then closes the problem entirely: setting in addition $\gamma_2=0$
(Model-$B_2$) collapses the characteristic family to vertical lines and degenerates the
PDE to an ODE in $x$ alone, whose boundary function $P_y(y)$ is pinned by an analyticity
argument at the interior singular point $x=y$ (Section~\ref{sec:model_b2}); removing
jockeying entirely while keeping only class-1 abandonment (Model-$C_2$) likewise
degenerates the PDE to an ODE in $x$, whose $P_y(y)$ is pinned by a boundedness argument
at $x=1$ (Section~\ref{sec:model_c2}); and zeroing all four parameters returns the fully
solvable baseline Model-$A$ (Section~\ref{sec:model_a}). The full model~$X$, retaining
all mechanisms, retains both obstructions simultaneously and is the one case that does
not close.
